\subsection{MONOGENOUS GROUPS}

Let \(a \in Z\) ; since \(aZ\) is a subgroup of \(Z\) , the relation between elements \(x, y\) of \(Z\) which states ``there exists \(z \in Z\) such that \(x - y = az\) '' is an equivalence relation, which we agree, once and for all, to write as \(x \equiv y \pmod{a}\) or \(x \equiv y(a)\) and which is called congruence modulo \(a\) . Replacing \(a\) by \(-a\) an equivalent relation is obtained, hence it may be supposed that \(a \geqslant 0\) ; for \(a = 0\) , \(x \equiv y(0)\) means \(x = y\) , hence a relation distinct from equality is obtained only if \(a \neq 0\) : we shall therefore suppose in what follows that \(a > 0\) unless otherwise indicated.

For \(a > 0\) , the quotient of \(\mathbf{Z}\) by the congruence \(x \equiv y(a)\) , that is the group \(\mathbf{Z}/a\mathbf{Z}\) , is called the group of rational integers modulo \(a\) .

PROPOSITION 16. Let \(a\) be an integer \(>0\) . The integers \(r\) such that \(0 \leqslant r < a\) form a system of representatives of the equivalence relation \(x \equiv y \pmod{a}\) on \(\mathbf{Z}\) .

If \(x\) is an integer \(\geqslant 0\) , there exist (Set Theory, III, § 5, no. 6) integers \(q\) and \(r\) such that \(x = aq + r\) and \(0 \leqslant r < a\) and \(x \equiv r \pmod{a}\) . If \(x\) is an integer \(\leqslant 0\) , the integer \(-x\) is \(\geqslant 0\) and by the above there exists an integer \(r\) such that \(0 \leqslant r < a\) and \(-x \equiv r \pmod{a}\) . Writing \(r' = 0\) if \(r = 0\) and \(r' = a - r\) if \(r > 0\) , then

\[
x \equiv - r \equiv r ^ {\prime} \pmod {a}
\]

and \(0 \leqslant r' < a\) . We now show that if \(0 \leqslant r < r' < a\) , then \(r \not\equiv r' \pmod{a}\) . Now \(r' - r < na\) for \(n \geqslant 1\) and \(r' - r > na\) for \(n \leqslant 0\) , whence \(r' - r \notin a\mathbf{Z}\) .

COROLLARY. Let \(a\) be an integer \(> 0\) . The group \(\mathbf{Z} / a\mathbf{Z}\) of rational integers modulo \(a\) is a group of order \(a\) .

PROPOSITION 17. Let H be a subgroup of Z. There exists one and only one integer \(a \geqslant 0\) such that H = aZ.

If \(\mathbf{H} = \{0\}\) , then \(\mathbf{H} = 0\mathbf{Z}\) . Suppose that \(\mathbf{H} \neq \{0\}\) . The subgroup \(\mathbf{H}\) has an element \(x \neq 0\) . Then \(x > 0\) or \(-x > 0\) , and therefore \(\mathbf{H}\) has elements \(>0\) . Let \(a\) be the smallest element \(>0\) in \(\mathbf{H}\) . The subgroup \(a\mathbf{Z}\) generated by \(a\) is contained in \(\mathbf{H}\) ; we show that \(\mathbf{H} \subset a\mathbf{Z}\) . Let \(y \in \mathbf{H}\) . By Proposition 16, there exists an integer \(r\) such that \(y \equiv r (\text{mod. } a)\) and \(0 \leqslant r < a\) . A fortiori \(y \equiv r (\text{mod. } H)\) , whence \(r \in \mathbf{H}\) . But this is only possible if \(r = 0\) and therefore \(y \in a\mathbf{Z}\) . The integer \(a\) is unique: if \(\mathbf{H} = \{0\}\) , then necessarily \(a = 0\) , and if \(\mathbf{H} \neq \{0\}\) , the integer \(a\) is the order of \(\mathbf{Z}/\mathbf{H}\) .

DEFINITION 15. A group is called monogenous if it admits a system of generators consisting of a single element. A finite monogenous group is called cyclic.

Every monogenous group is commutative (no. 3, Corollary 2 to Proposition 2). Every quotient group of a monogenous group is monogenous (no. 3, Corollary 3 to Proposition 2).

The additive group Z is monogenous: it is generated by \(\{1\}\) . For every positive integer a, the group Z/aZ is monogenous, for it is a quotient of Z.

PROPOSITION 18. A finite monogenous group of order \(a\) is isomorphic to \(\mathbf{Z}/a\mathbf{Z}\) . An infinite monogenous group is isomorphic to \(\mathbf{Z}\) .

Let G be a monogenous group (written multiplicatively) and x a generator of G. The identity \(x^{m}x^{n} = x^{m+n}\) (§ 1, no. 3, formula (1)) shows that the mapping \(n \mapsto x^{n}\) is a homomorphism of Z into G. Its image is a subgroup of G containing x and hence it is G. By no. 5, Theorem 3, the group G is isomorphic to the quotient of Z by a subgroup, which is necessarily of the form aZ (Proposition 17). If a \textgreater{} 0, the group G is finite of order a and if a = 0, the group G is isomorphic to Z.

PROPOSITION 19. Let \(a\) be an integer \(>0\) . Let \(\mathbf{H}\) be a subgroup of \(\mathbf{Z}/a\mathbf{Z}\) , \(b\) the order of \(\mathbf{H}\) and \(c\) its index in \(\mathbf{Z}/a\mathbf{Z}\) . Then \(a = bc\) , \(\mathbf{H} = c\mathbf{Z}/a\mathbf{Z}\) and \(\mathbf{H}\) is isomorphic to \(\mathbf{Z}/b\mathbf{Z}\) .

Conversely, let \(b\) and \(c\) be two integers \(>0\) such that \(a = bc\) . Then \(a\mathbf{Z} \subset c\mathbf{Z}\) and \(c\mathbf{Z}/a\mathbf{Z}\) is a subgroup of \(\mathbf{Z}/a\mathbf{Z}\) , of order \(b\) and index \(c\) .

a = bc by no. 4, Corollary to Proposition 4. By no. 7, Theorem 4, H is of the form \(H'/aZ\) , where \(H'\) is a subgroup of Z and \(Z/H'\) is isomorphic to \((\mathbf{Z}/a\mathbf{Z})/H\) and hence of order c.~By Proposition 17 and the Corollary to Proposition 16, \(H' = cZ\) and hence H is monogenous. Finally, H is isomorphic to Z/bZ by Proposition 18. Conversely, if a = bc, then \(aZ \subset cZ\) for \(a \in cZ\) : the quotient group \((\mathbf{Z}/a\mathbf{Z})/(c\mathbf{Z}/a\mathbf{Z})\) is isomorphic to Z/cZ (no. 7, Theorem 4) and hence of order c (no. 4, Corollary to Proposition 4) and index b (no. 4, Corollary to Proposition 4).

COROLLARY. Every subgroup of a monogenous group is monogenous.

Let \(a\) and \(b\) be two integers \(\neq 0\) . The relation \(b \in a\mathbf{Z}\) is also written: \(b\) is a multiple of \(a\) , and also \(a\) divides \(b\) or \(a\) is a divisor of \(b\) .

DEFINITION 16. An integer \(p > 0\) is called prime if \(p \neq 1\) and it admits no divisor \(> 1\) other than \(p\) .

PROPOSITION 20. An integer \(p > 0\) is prime if and only if the group \(\mathbf{Z} / p\mathbf{Z}\) is a simple group.

This follows from Proposition 19.

COROLLARY. Every commutative simple group is cyclic of prime order.

Let G be such a group. Then \(G \neq \{e\}\) ; let \(a \neq e\) be an element of G. The subgroup generated by a is normal since G is commutative, it is not reduced to \(\{e\}\) and hence is equal to G. Therefore G is monogenous and hence isomorphic to a group of the form Z/pZ with p \textgreater{} 0, since Z is not simple, and p is necessarily prime.

Remark. A finite group G of prime order is necessarily cyclic. G admits no subgroup other than G and \(\{e\}\) and hence it is generated by every element \(\neq e\) .

Lemma 2. Let a be an integer \textgreater0. By associating with every composition series \((\mathbf{H}_{i})_{0\leqslant i\leqslant n}\) of the group Z/aZ the sequence \((s_{i})_{1\leqslant i\leqslant n}\) , where \(s_{i}\) is the order of \(H_{i-1}/H_{i}\) , a one-to-one correspondence is obtained between the composition series of Z/aZ and the finite sequences \((s_{i})\) of integers \textgreater0 such that \(a=s_{1}\ldots s_{n}\) . The composition series \((\mathbf{H}_{i})_{0\leqslant i\leqslant n}\) is a Jordan-Hölder series if and only if the \(s_{i}\) are prime.

If \((\mathbf{H}_i)_{0\leqslant i\leqslant n}\) is a composition series of \(\mathbf{Z} / a\mathbf{Z}\) , it follows, by induction on \(n\) , from no. 4, Proposition 4, that \(a = \prod_{i=1}^{n} (\mathbf{H}_{i-1}:\mathbf{H}_i)\) .

Conversely, let \((s_i)_{1 \leqslant i \leqslant n}\) be a sequence of integers \(>0\) such that \(a = s_1 \ldots s_n\) . If \((\mathbf{H}_i)_{0 \leqslant i \leqslant n}\) is a composition series of \(\mathbf{Z} / a\mathbf{Z}\) such that \((\mathbf{H}_{i-1} : \mathbf{H}_i) = s_i\) for \(1 \leqslant i \leqslant n\) , then necessarily \(((\mathbf{Z} / a\mathbf{Z}) : \mathbf{H}_i) = \prod_{1 \leqslant j \leqslant i} s_j\) as is seen by induction on \(i\) , whence \(\mathbf{H}_i = \left( \prod_{j=1}^{i} s_j \right) \mathbf{Z} / a\mathbf{Z}\) (Proposition 19), which shows the injectivity of the mapping in question. We now show its surjectivity. Writing \(H_{i}=\left(\prod_{j=1}^{i}s_{j}\right)\mathbf{Z}/a\mathbf{Z}\) for \(0\leqslant i\leqslant n\) , a composition series of Z/aZ is obtained such that \((\mathrm{H}_{i-1}:\mathrm{H}_{i})=s_{i}\) (Proposition 19). The second assertion follows from Proposition 20 and no. 7, Proposition 9.

Let \(\mathfrak{P}\) denote the set of prime numbers.

THEOREM 7 (decomposition into prime factors). Let \(a\) be a strictly positive integer. There exists one and only one family \((v_{p}(a))_{p \in \mathfrak{P}}\) of integers \(\geqslant 0\) such that the set of \(p \in \mathfrak{P}\) with \(v_{p}(a) \neq 0\) is finite and

\[
a = \prod_ {p \in \mathfrak {P}} p ^ {v _ {p} (a)}.\tag{5}
\]

As the group \(\mathbf{Z}/a\mathbf{Z}\) is finite, it admits a Jordan-Hölder series. Lemma 2 then implies that \(a\) is a product of prime integers, whence the existence of the family \((v_{p}(a))\) ; further, for every family \((v_{p}(a))_{p\in\mathfrak{P}}\) satisfying the conditions of Theorem 7, the integer \(v_{p}(a)\) is, for all \(p\in\mathfrak{P}\) , equal to the number of factors of a Jordan-Hölder series of \(\mathbf{Z}/a\mathbf{Z}\) isomorphic to \(\mathbf{Z}/p\mathbf{Z}\) (Lemma 2). The uniqueness of the family \((v_{p}(a))\) therefore follows from the Jordan-Hölder theorem (no. 7, Theorem 6).

COROLLARY. Let \(a\) and \(b\) be two integers \(>0\) . Then \(v_{p}(ab) = v_{p}(a) + v_{p}(b)\) . For \(a\) to divide \(b\) , it is necessary and sufficient that \(v_{p}(a) \leqslant v_{p}(b)\) for every prime number \(p\) .

In any group G, if the (monogenous) subgroup generated by an element \(x \in G\) is of finite order d, x is called an element of order d; the number d is therefore the least integer \textgreater0 such that \(x^{d} = e\) ; if the subgroup generated by x is infinite, x is said to be of infinite order. These definitions, together with Proposition 4 (no. 4), imply in particular that in a finite group G the order of every element of G is a divisor of the order of G.

PROPOSITION 21. In a finite group G of order n, \(x^{n} = e\) for all \(x \in G\) .

If \(p\) is the order of \(x\) , then \(n = pq\) , with \(q\) an integer, and hence \(x^n = (x^p)^q = e\) .

